%******************************************************************************%
% Copyright (C) 2018  Louis Solofrizzo                                         %
%                                                                              %
% This content is considered a free software: you can redistribute it          %
% and/or modify it under the terms of the GNU General Public License as        %
% published by the Free Software Foundation, either version 3 of the License,  %
% or (at your option) any later version.                                       %
%                                                                              %
% This program is distributed in the hope that it will be useful,              %
% but WITHOUT ANY WARRANTY; without even the implied warranty of               %
% MERCHANTABILITY or FITNESS FOR A PARTICULAR PURPOSE.  See the                %
% GNU General Public License for more details.                                 %
%                                                                              %
% You should have received a copy of the GNU General Public License            %
% along with this program.  If not, see <https://www.gnu.org/licenses/>.       %
%******************************************************************************%

%******************************************************************************%
%                                                                              %
%                          KFS_9.en.tex for KFS_9                              %
%                                                                              %
%                  Created on : Wed May 25 13:27:28 2016                       %
%          Made by : Louis "Ne02ptzero" Solofrizzo <louis@ne02ptzero.me>       %
%                                                                              %
%******************************************************************************%

\documentclass{42-en}


%******************************************************************************%
%                                                                              %
%                                    Header                                    %
%                                                                              %
%******************************************************************************%
\begin{document}



                           \title{KFS\_9}
                          \subtitle{ELF}
                       \member{Louis Solofrizzo}{louis@ne02ptzero.me}
                        \member{42 Staff}{pedago@42.fr}

\summary {
    Finally some execution.
}

\maketitle

\tableofcontents


%******************************************************************************%
%                                                                              %
%                                  Foreword                                    %
%                                                                              %
%******************************************************************************%
\chapter{Foreword}

%******************************************************************************%
%                                                                              %
%                                 Introduction                                 %
%                                                                              %
%******************************************************************************%
\chapter{Introduction}
    \subsection{ELF}
        \begin{quotation}
        \textit{In computing, the Executable and Linkable Format (ELF, formerly
        called Extensible Linking Format) is a common standard file format for
        executables, object code, shared libraries, and core dumps. First
        published in the System V Release 4 (SVR4) Application Binary Interface
        (ABI) specification, and later in the Tool Interface Standard, it
        was quickly accepted among different vendors of Unix systems. In 1999
        it was chosen as the standard binary file format for Unix and Unix-like
        systems on x86 by the 86open project.
        ELF is flexible and extensible by design, and it is not bound to any
        particular processor or architecture. This has allowed it to be adopted
        by many different operating systems on many different platforms.}
        \end{quotation}
        As you can read above, ELF is a format used to read binary files.
        It is pretty common in many OSs, and now's the time to implement it!\\
        More info \href{http://wiki.osdev.org/ELF}{Here.}

%******************************************************************************%
%                                                                              %
%                                  Goals                                       %
%                                                                              %
%******************************************************************************%
\chapter{Goals}
    At the end of this subject, you will have:
    \begin{itemize}\itemsep1pt
        \item A complete interface to read, parse, stock and execute ELF files.
        \item Syscalls to read ELF files and launch a process with them.
        \item A kernel module in ELF, ready to be inserted at run time.
    \end{itemize}

%******************************************************************************%
%                                                                              %
%                             General instructions                             %
%                                                                              %
%******************************************************************************%
\chapter{General instructions}
    \section{Code and Execution}
        \subsection{Emulation}
        The following part is not mandatory, you're free to use any virtual
        manager you want; however, I suggest you use \texttt{KVM}.
        It's a \texttt{Kernel Virtual Manager} with advanced execution
        and debug functions.
        All of the examples below will use \texttt{KVM}.
        \subsection{Language}
            The \texttt{C} language is not mandatory, you can use any language
            you want for this series of projects.\\
            Keep in mind that not all languages are kernel friendly though, you
            could code a kernel in \texttt{Javascript}, but are you sure it's a
            good idea?\\
            Also, most of the documentation is written in \texttt{C}, you will
            have to 'translate' the code all along if you choose a different
            language.\\

            Furthermore, not all the features of a given language can be used
            in a basic kernel. Let's take an example with \texttt{C++}:\\
            this language uses 'new' to make allocations, classes and
            structures declarations. But in your kernel you don't have a memory
            interface (yet), so you can't use any of these features.\\

            Many languages can be used instead of \texttt{C},
            like \texttt{C++}, \texttt{Rust}, \texttt{Go}, etc.
            You can even code your entire kernel in \texttt{ASM}!\\
            So yes, you may choose a language. But choose wisely.
            \begin{center}
              \includegraphics[width=8cm]{choose.jpg}
            \end{center}
\newpage

    \section{Compilation}
        \subsection{Compilers}
            You can choose any compiler you want. I personaly use \texttt{gcc}
            and \texttt{nasm}. A Makefile must be turned-in as well.
        \subsection{Flags}
            In order to boot your kernel without any dependency, you must
            compile your code with the following flags (adapt the flags for
            your language, these are \texttt{C++} examples):
            \begin{itemize}\itemsep1pt
                \item \texttt{-fno-builtin}
                \item \texttt{-fno-exception}
                \item \texttt{-fno-stack-protector}
                \item \texttt{-fno-rtti}
                \item \texttt{-nostdlib}
                \item \texttt{-nodefaultlibs}
            \end{itemize}
            You might have noticed these two flags: \texttt{-nodefaultlibs}
            and \texttt{-nostdlib}. Your Kernel will be compiled on a host
            system, that's true, but it cannot be linked to any existing
            library on that host, otherwise it will not be executed.
    \section{Linking}
        You cannot use an existing linker in order to link your kernel.
        As mentionned above, your kernel would not be initialized. So you must
        create a linker for your kernel.\\
        Be careful, you \texttt{CAN} use the 'ld' binary available on your
        host, but you \texttt{CANNOT} use the .ld file of your host.
    \section{Architecture}
        The \texttt{i386} (x86) architecture is mandatory (you can thank
        me later).
    \section{Documentation}
        There is a lot of documentation available, good and bad.
        I personaly think the \texttt{\href{http://wiki.osdev.org/Main_Page}
        {OSDev}} wiki is one of the best.
    \section{Base code}
        In this subject, you have to take your previous \texttt{KFS} code,
        and work from it!\\
        Or don't. And rewrite everything from scratch. Your call!
\newpage

%******************************************************************************%
%                                                                              %
%                             Mandatory part                                   %
%                                                                              %
%******************************************************************************%
\chapter{Mandatory part}
    In this project, you will have to:
    \begin{itemize}\itemsep1pt
        \item Create a parser for ELF files.
        \item Create a loader for ELF files.
        \item Execute an ELF file.
        \item Implement these points in a \texttt{execve (2)}-like
        syscall.
        \item Make the syscall create a process associated with the ELF file.
        \item Rewrite your API modules, so that you can load
        modules in ELF at run time.
        \item Load the modules, list them and unload them.
    \end{itemize}

%******************************************************************************%
%                                                                              %
%                                 Bonus part                                   %
%                                                                              %
%******************************************************************************%
\chapter{Bonus part}
    You must create a memory ring for the modules (specifically, built-in and
    run-time modules).

%******************************************************************************
%                                                                              %
%                           Turn-in and peer-evaluation                        %
%                                                                              %
%******************************************************************************%
\chapter{Turn-in and peer-evaluation}

    Turn your work in using your \texttt{GiT} repository, as
    usual. Only the work that's in your repository will be graded during
    the evaluation.

    Your must turn in your code, a Makefile and a basic virtual image for your
    kernel.\\
    Side note about this image,
    THERE IS NO NEED TO BE BUILT LIKE AN ELEPHANT.

%******************************************************************************%
\end{document}
